% chapters/07_conclusion.tex
% Rewritten 2026-07-21 from a claim-and-scope matrix for Chapters 4--6.
% Objective: conclusion-level synthesis rather than another chapter recap. The
% three studies are treated as distinct evidential cases, not points on a common
% quantitative scale. Information theory describes dependence; geometry describes
% organization and accessibility. They are complementary, not interchangeable.
% Framing locks: guaranteed failure is Ch.4-only; Ch.5 is diagnostic and recoverable;
% Ch.6 claims parity, reports the Natural Questions losses, and never claims SOTA.
% No new results, re-derivations, or hard-coded result values appear here.
\chapter{Conclusion}
\label{ch:conclusion}

A reliability detector can fail for two different reasons. Its scoring rule may
be too weak to extract a distinction that the model represents, or the model may
not represent that distinction in the first place. These cases call for different
responses. A better readout can address the first. The second requires a change to
the learned features or the objective used to train them. The detector may instead
need access to additional information. This dissertation has separated those cases
across two forms of out-of-distribution (OOD) detection and one form of
hallucination detection.

Information theory and representation geometry made that separation possible.
Information theory asks whether the representation remains statistically dependent
on the target decision. Geometry describes how relevant variation is arranged
across feature directions and whether a particular detector can use it.
The two perspectives address the same representational question without measuring
an identical object. The three studies likewise form a conceptual progression,
not a shared numerical scale: \Cref{ch:label-blindness} establishes a conditional
failure boundary, \Cref{ch:domain-sensitivity-collapse} diagnoses and repairs
suppressed sensitivity, and \Cref{ch:hallucination-detection} constructs a detector
for a signal that remains accessible inside a frozen model. This chapter states
what those results jointly establish, where the common framework stops, and which
questions follow from it.

\section{What the Dissertation Establishes}
\label{sec:concl:summary}

Each study supplies a different kind of evidence. The first proves a limit under a
stated independence condition. The second gives conditional bounds and tests a
training-time remedy. The third uses a feasibility argument to motivate a detector
and evaluates the resulting construction. Their conclusions should therefore be
combined without treating their guarantees as interchangeable.

\subsection{A Conditional Limit on Unlabeled OOD Detection}

\Cref{ch:label-blindness} asks whether OOD detection can use features produced
without the labels that define the eventual task. When the self-supervised or
unsupervised surrogate task is independent of the label-relevant features, the
minimal sufficient representation for that task contains no
information about the required label distinction. A detector restricted to that
representation then has no statistic from which to recover the distinction. The
Label Blindness Theorem makes this a \emph{guaranteed-failure} result under its
independence condition (\Cref{mainbodyfailood}). It is not a claim that every
unlabeled representation fails, nor that strong performance on conventional OOD
benchmarks is impossible.

The Adjacent OOD benchmark makes the condition empirically visible. Many benchmark
pairs differ through dataset-level cues that an unlabeled representation can use
without encoding the held-out class distinction. Adjacent OOD increases the overlap
between in-distribution and OOD inputs so that incidental cues provide less help.
High overlap does not by itself imply label blindness; the independence condition
remains essential. The benchmark instead tests the setting in which dependence on
the relevant label features matters most. Together, the theorem and benchmark
distinguish success on a chosen dataset pair from a general guarantee that an
unlabeled objective will preserve whatever a future OOD task requires.

\subsection{A Geometric Diagnosis and Training-Time Recovery}

\Cref{ch:domain-sensitivity-collapse} studies a different failure. Here labels are
available, but single-domain supervision rewards separation among the known classes
without directly rewarding sensitivity to changes of domain. Domain-Sensitivity
Collapse (DSC) describes the resulting concentration of variance in a
class-discriminative subspace and the suppression of directions that would respond
to domain shift. Under the conditions developed in the chapter, distance-based and
logit-based OOD scores lose separation because both operate on the affected
features (\Cref{thm:dsc:dist-fail,prop:dsc:msp-insens}). These bounds diagnose a
mechanism. They do not prove that every single-domain model must fail or that no
intervention can recover the signal.

Teacher-Guided Training (TGT) tests the diagnosis by acting on the representation
rather than replacing the scorer. It transfers the class-suppressed residual of a
frozen multi-domain teacher into the student during training, after which the
teacher and auxiliary head are discarded (\Cref{sec:dsc:method}). The resulting
changes in feature geometry accompany improvements for the post-hoc scorers studied
in the chapter, with the most consistent gains in the distance-based family that
the DSC analysis directly addresses. TGT therefore supplies more than a second OOD
method: it shows that a geometric account of lost sensitivity can identify a
training-time intervention, while retaining ordinary student-only inference.

\subsection{Detection from Cross-Layer Language-Model Activations}

\Cref{ch:hallucination-detection} begins from a frozen language model, where
retraining the underlying representation is outside the detection setting. The
question is whether intermediate activations contain a usable distinction between
faithful and hallucinated generations. Dependence between layers makes a
contrastive objective over layer-pair views well-posed. The chapter uses this
feasibility argument to motivate Contrastive+Recon, a one-class contrastive probe
with an auxiliary reconstruction channel (\Cref{sec:hd:method}). Cross-layer
dependence alone does not establish that hallucination labels are recoverable.
Evidence comes from the detector's discrimination performance and the predicted
failure of symmetric supervision (\Cref{sec:hd:results,sec:hd:ablations}).

Across the main comparison, Contrastive+Recon occupies the same performance band
as the strongest engineered probe. It \emph{matches-or-outperforms, in the mean,}
that baseline. This is a parity claim rather than dominance. The method wins some
comparisons, ties others within the observed variation, and loses on Natural
Questions for both models. Its cross-dataset results also support parity under
transfer rather than a transfer win (\Cref{sec:hd:transfer}). These outcomes support the cross-layer design
rationale. They do not show that the probe is an optimal readout or that the
information-theoretic argument selects the best detector for every dataset.

\subsection{A Joint Finding About Representation and Detection}
\label{subsec:concl:synthesis}

The three studies answer one question under different constraints: does the model
representation available to the detector contain the distinction the detector must
make? \Cref{ch:label-blindness} studies a condition under which the answer is no.
\Cref{ch:domain-sensitivity-collapse} studies features whose domain sensitivity has
been suppressed by training and shows that the geometry can be changed.
\Cref{ch:hallucination-detection} studies a frozen model for which the relevant
signal remains accessible in intermediate activations. The failure, recovery, and
detection sequence names these cases. It does not rank the studies by a common
amount of information because the target and evidence differ across models.

What carries across them is a constraint on method design. Improving a scorer
cannot recover a distinction that is absent from every input the scorer receives.
When training has reduced sensitivity to a distinction, changing the representation
may be more effective than continuing to refine a post-hoc score. When the signal
remains available, the readout and its supervision determine how much of that signal
is usable in practice. Representation content is therefore a necessary part of a
detection analysis, but it is not sufficient by itself. Accessibility and
estimation still matter, as does detector capacity.

This distinction is the dissertation's common contribution. It changes the first
question asked about a failed detector. Instead of immediately selecting another
score, the analysis asks what distinction the score can observe and whether training
preserved it. The available evidence then determines whether to change the readout
or the representation. The three studies do not produce a universal detector. They
provide three ways to answer that prior question and demonstrate the different
interventions that follow from its answer.

\section{Implications of the Common Framework}
\label{sec:concl:implications}

The common framework matters beyond the individual methods because it separates
problems that are often combined under the broad label of detection failure. It
also clarifies what the two analytical perspectives contribute and how evaluation
can test a proposed mechanism.

\subsection{Representation Failure and Detector Failure}
\label{subsec:concl:representation}

An observed detection error does not reveal whether the representation or the
scoring rule caused it. Benchmarking several scores on fixed features can compare
readouts, but it cannot show that the features contain every distinction the task
requires. Label blindness supplies the clearest case: once the stated independence
condition removes the target distinction, no post-hoc score on that representation
can reconstruct it. DSC supplies a less absolute case in which existing features
respond weakly to domain shift and training can restore useful directions. The
hallucination study addresses the remaining case, where a single-pass readout can
extract a distinction from activations already computed by the model.

These cases suggest an intervention order. First determine what the detector can
observe. If the required signal is absent, the learning objective or representation
must change unless the detector receives new inputs. External retrieval and repeated
generations can change what is observable without retraining the base model. If the
signal has been suppressed but remains recoverable, a
training-time objective can strengthen it. If it is already available, readout
design becomes the relevant problem. This order is a practical diagnostic, not a
theorem that guarantees the chosen intervention will succeed.

\subsection{Information Theory and Geometry as Complementary Evidence}
\label{subsec:concl:lenses}

Information theory and representation geometry contribute different evidence to
this diagnostic. Mutual information states whether statistical dependence remains
between a representation and a target. It can establish an impossibility boundary
when that dependence vanishes, or justify an objective when a lower bound can be
optimized. It does not, by itself, show that a simple detector can recover the
dependence from finite data. Geometry describes where variation lies, how strongly
it is concentrated, and whether the directions used by a scorer respond to the
shift. A favorable geometric arrangement can make a signal accessible, but a
geometric statistic alone does not establish its relationship to every target.

The studies use the perspective suited to each claim. Label blindness and the
cross-layer probe are expressed through dependence and information bounds. DSC uses
feature spectra and subspaces, then relates them to scorer sensitivity. Similar training
processes may affect both information and geometry, but this dissertation does not
prove that mutual-information loss and rank collapse are equivalent. Its more
defensible methodological conclusion is that reliability analysis benefits from
both questions: whether target dependence remains and how the surviving variation
is organized for the detector that must use it.

\subsection{Evaluation as a Test of Mechanism}
\label{subsec:concl:eval}

The experiments also illustrate a role for evaluation beyond producing an aggregate
ranking. A theoretical or mechanistic claim should identify a regime in which its
predicted effect becomes visible. Adjacent OOD reduces incidental separation so
that label dependence becomes consequential. The single-domain experiments measure
feature rank alongside OOD performance and compare in-domain with out-of-domain
shifts, allowing the geometric diagnosis and the proposed recovery to be examined
together. In hallucination detection, compute-matched comparisons separate the
benefit of activation access from the cost of repeated sampling. The symmetric
supervision ablation and cross-dataset evaluation test predictions about the learned
signal rather than only its average score.

Negative and uneven results are informative under this view. The Natural Questions
losses identify a setting in which the cross-layer rationale does not match the
strongest engineered probe. Variable TGT gains identify where teacher--student
complementarity or scorer choice affects recovery. Strong performance by unlabeled
methods on far-OOD pairs shows why label blindness may be hidden rather than
disproving the conditional result. A benchmark supports a reliability claim only
to the extent that its design exposes the mechanism and scope named by that claim.

\section{Scope and Limitations}
\label{sec:concl:limitations}

The common framework does not amount to a universal necessary-and-sufficient theory
of detectability. It combines one conditional impossibility result, one conditional
geometric diagnosis with an empirical remedy, and one supported feasibility
argument. None establishes that representation quality alone determines detection
performance, and none identifies an optimal detector for every distribution.

\paragraph{Label blindness.}
The strict result assumes independence between the surrogate task and the
label-relevant features. Real self-supervised objectives often share some structure
with later semantic tasks, so the exact theorem describes a boundary rather than an
expected field failure rate. Approximate label blindness is correspondingly harder
to diagnose because small dependence may still be useful to a strong readout. The
Adjacent OOD construction examines substantial feature overlap, which makes the
failure easier to expose. It does not represent the full range of shifts encountered in
deployment. Overlap and independence must not be treated as the same condition.

\paragraph{Domain-Sensitivity Collapse.}
The DSC bounds apply under stated assumptions about feature anisotropy and dominant-
subspace matching, together with limits on scorer sensitivity. They explain when the
studied scores lose separation; they are not convergence results across every
single-domain training setting.
The empirical evaluation covers a varied but finite collection of domains and shows
that improvements differ across scorers and teacher--student pairings. TGT also
depends on the teacher containing domain-sensitive structure that the student lacks.
A teacher with limited coverage or a representation too similar to the student can
provide a weaker recovery signal.

\paragraph{Hallucination detection.}
Contrastive+Recon and the strongest engineered baseline occupy the same performance
band rather than separating consistently. The Natural Questions losses remain
unexplained, and the transfer results do not remove them. The feasibility argument
supports the use
of cross-layer views but does not prove that the selected layers or one-class
objective extract the available signal optimally. The same uncertainty applies to
the reconstruction channel and distance scorer. The evaluation is also limited to
the studied open-weights models and short-answer tasks. It requires white-box
activation access and uses operational labels produced by an automatic evaluator.
These constraints limit both applicability and the precision of the target being
learned (\Cref{sec:hd:limitations}).

\paragraph{The synthesis.}
The absent, suppressed, and accessible cases organize the contributions, but they
are not exhaustive states with fixed thresholds. A representation may preserve a
weak signal that is usable only with more data or a nonlinear readout. Additional
context may also be necessary. The detector can introduce information through
retrieval or sampling, and another external source can change the task again. The
framework should therefore separate sources of failure and select the next analysis;
it is not a complete scalar measure of reliability.

\section{Future Directions}
\label{sec:concl:future}

The limitations lead to research questions that test and extend the framework
rather than simply applying the existing methods to more benchmarks.

\paragraph{Measure label blindness by degree.}
A useful diagnostic would estimate how much dependence a particular learning
objective preserves about a proposed label distinction. Such a measure would need
to account for the data distribution as well as the objective and labels; the risk
cannot be inferred from the objective name alone. Estimation before or during
training could turn strict label blindness from a boundary result into a screening
tool, while experiments could determine how much residual dependence different
detectors require.

\paragraph{Relate information loss to geometric collapse under explicit
conditions.}
DSC was developed geometrically, while the other studies rely primarily on
information-theoretic arguments. Future work could identify assumptions under which
changes in target mutual information imply changes in the feature spectrum or in a
shift-sensitive subspace, and also construct counterexamples where they diverge.
This would test the relationship between the two perspectives without presuming
that they measure the same object. It could also provide diagnostics that connect a
geometric warning to the target dependence a deployment actually requires.

\paragraph{Recover domain sensitivity beyond a fixed teacher.}
TGT inherits both the strengths and the blind spots of its multi-domain teacher.
Future methods could compare teachers with different coverage and select residual
targets adaptively. Another approach would recover domain-sensitive variation
without a fixed teacher.
The unresolved in-domain OOD regime is especially important because class and domain
cues overlap there more strongly. Progress should be evaluated through changes in
the representation as well as changes in the final OOD score, so that an improved
number can be tied to the proposed recovery mechanism.

\paragraph{Generalize and act on the hallucination signal.}
The cross-layer construction should be tested across model families and scales,
including different generation formats, with particular attention to the Natural
Questions shortfall and to settings where activations are unavailable. Layer
selection remains open, as do the relative roles of contrastive and reconstruction
supervision. A further
step would move from detection to mitigation by asking whether the identified
directions can safely alter generation rather than merely flag it. That question
requires a causal evaluation of the intervention and cannot be answered by detector
performance alone.

\section{Concluding Remarks}
\label{sec:concl:remarks}

This dissertation began from a practical problem: models can remain confident when
their inputs or outputs fall outside what their training supports. Its central
conclusion is not that one detector solves this problem across domains. It is that
detector design should begin by establishing what distinction the available
representation contains. \Cref{ch:label-blindness} marks a condition under which
that distinction is absent. \Cref{ch:domain-sensitivity-collapse} shows that
training can restore suppressed sensitivity, while
\Cref{ch:hallucination-detection} demonstrates a readout for a signal retained
inside a frozen model.

The resulting framework separates limits from repairable failures and repairable
failures from readout problems. Information theory and geometry provide different
evidence for making those distinctions. Their joint value lies in directing the
analysis to the representation before another score is added on top.
When the representation keeps it, a detector can read it out.
